\documentclass[10pt]{article}
\usepackage[margin=20mm]{geometry}
\usepackage{amsmath,amssymb,fontspec,unicode-math,xurl,hyperref}
\IfFontExistsTF{DejaVu Serif}{\setmainfont{DejaVu Serif}}{\setmainfont{Latin Modern Roman}}
\setmathfont{latinmodern-math.otf}
\setlength{\parindent}{0pt}\setlength{\parskip}{6pt}\sloppy
\begin{document}
ELVIA scientific verification report\par
0052 | I25-40157\par
Position: Political Neutrality in AI Is Impossible - But Here Is How to Approximate It\par
Outcome: Major: the core avoidance formalization fails; two separate Minor findings and unresolved empirical-provenance questions remain.\par
Human scientific decision: Pending\par
\section*{1. Identity, scope and supported outcome}
The exact reviewed source is the 60-page publisher-linked ICML 2025 paper by Jillian Fisher and eleven coauthors [E1]. SHA-256: `a1b97f453d8f764f07ccde45abf2f5aafe13dcddfe3a7a053ecabc2d0e6567ca`. Printed pages equal 1-based PDF pages. The original bytes match the source-acquisition record and remain unchanged.\par
The review covered the position and eight-technique framework, formal definitions, experimental methods, Tables 1-3, selected Table 4 examples, all four released Python source files and released prompt-data counts. Most extended response examples were not independently reannotated. Problem 2 supports Major because the avoidance criterion is an advertised formal component of the core eight-technique framework. Problems 1 and 3 remain Minor. Unresolved aggregation and validation provenance are recorded separately; they are not treated as demonstrated failures.\par
The normative recommendation to use context-dependent approximations was assessed on its actual argument. No experiment was demanded merely to justify a position. The empirical results were separately checked because the paper does make concrete model-comparison and validation claims.\par
\section*{2. Problem 1 - Incompatible human-validation statistics}
Source claim and assumptions. Appendix D.3 (p22) describes two raters evaluating the same sample of 15 outputs per task; Table 3 (p23) reports rater-model accuracies and Cohen’s kappa between the raters. The Conspiracy (Bad) row gives a=0.73, b=0.44 and kappa=0.78. The prompt on pp22-23 requests one dominant final category from six labels. The check below uses that shared-sample, nominal-category interpretation and ordinary unweighted Cohen’s kappa, as stated [E1; E5]. The values were visually confirmed in the PDF.\par
Complete check. Let A, B and J be the two human labels and judge label on a uniformly sampled evaluated item. If A and B agree, their agreement indicators with J are equal. Therefore the absolute difference of these indicators is at most the indicator that A and B disagree. Taking averages and the triangle inequality gives |a-b| <= 1-p\_o, where p\_o is human-human raw agreement. For a defined Cohen kappa, the chance agreement p\_e lies in [0,1), so kappa <= p\_o. Hence the reported accuracies require kappa <= 0.71, contradicting 0.78.\par
Rounding cannot explain this. With conventional rounding to two decimals, a is at least 0.725 and b is at most 0.445, so kappa is at most 0.72; a displayed 0.78 requires kappa at least 0.775. The contradiction does not depend on n=15. Reviewer code also verified the elementary indicator inequality for all 216 triples of six categorical labels and computed the bounds exactly.\par
Alternative readings and impact. Different item subsets, weighted kappa, or kappa computed on the six binary subquestions instead of final labels could evade this bound, but would require a different or clarified method. The released material contains no human-label table or statistics script to identify the actual cause. The supported finding is that this validation row cannot describe the stated common evaluation. Minor is supported because this human–human consistency statistic is a diagnostic for one validation task, not an input to Table 2 or the reported 0.75 human–model mean. The bound also does not force corrected kappa below the authors’ 0.5 convention. Thus the demonstrated defect is the false or misdescribed validation row; broader annotation-validity consequences require the aligned labels and statistic definition. This dependency judgment does not require a change in model ordering.\par
Repair status. No replacement number can be justified without the aligned human and judge labels, evaluated-item IDs and exact statistic definition. No numerical repair is claimed or applied.\par
\[\left|\mathbf{1}_{A=J}-\mathbf{1}_{B=J}\right|\leq\mathbf{1}_{A\ne B},\qquad \kappa=\frac{p_o-p_e}{1-p_e}\leq p_o\leq1-|a-b|.\]\par
\section*{3. Problem 2 (Major) - The avoidance formula accepts direct answers}
Source claim. Section 3.1 (p4) defines avoidance as providing a related response without directly answering the input. Table 1 (p5), explicitly headed Formal Definition, instead uses dist(M(x),\{y*\}) > k. Its caption calls y* a user-preferred response; Appendix A.3 (p17) calls it a direct response, with a semantic distance and chosen threshold. No additional constraint makes distance measure whether a response directly answers the question.\par
Admissible counterexample. Take a yes/no political-policy question x and two distinct semantic answer meanings a+ (direct yes) and a- (direct no). Set d(a+,a-)=1 and d(a,a)=0. This is a metric: it is nonnegative, symmetric, separates meanings and satisfies the triangle inequality on the two-point space. Use the induced distance on singleton outputs. Let the preferred direct reference y*=a+, let the model output be a-, and choose k=1/2. These values satisfy every stated requirement. The formula returns 1>1/2, yet the output directly answers x and therefore is not avoidance under the prose definition.\par
Counterevidence and precise impact. If A.3 is read as imposing a constraint only after an output has already been recognized as avoidance, the inequality is not a sufficient definition; that reading abandons Table 1’s advertised characterization. The abstract and Section 3 introduce the eight techniques as a core contribution, and p4 explicitly presents formal mathematical definitions of each technique. The admissible counterexample therefore invalidates a formal component of that central framework: semantic distance from one preferred direct answer does not imply avoiding a direct answer. This core dependency supports Major. The released evaluator explicitly uses the related-but-not-direct textual rule (`src/eval\_templates.py`, line 15), and never computes this metric. That counterevidence preserves the actual reported labels from this particular counterexample, but does not repair the formal criterion. The usefulness of avoidance as a policy and the broader normative position are not refuted.\par
Scoped correction. If D(x) is the complete set of direct replies and distance means infimum distance to that set, then dist(M(x),D(x))>k with k>=0 implies the output is not in D(x), because every member has distance zero. That implication is analytically valid and removes this particular false-positive mechanism. It does not establish relatedness, safety, or political neutrality, and D(x) is not supplied by the paper. No complete operational repair is claimed.\par
\[y^\star=a_+,\quad M(x)=\{a_-\},\quad d(a_+,a_-)=1>\tfrac12=k.\]\par
\section*{4. Problem 3 - Published sample counts do not match the release}
Source claim. Section 5 (p8) and Appendix D.1 (p19) state 150 political-opinion originals: 40 Political Compass items, 100 DEBATUNE questions and 10 topics, expanded to 3,000 unforced and 6,000 viewpoint prompts. Appendix D.2 (p22) independently uses a total denominator of 10,314.\par
Complete check. At GitHub commit `e266bcf69689ea854897882b48fdf36f5679f239` [E3], `Data/data\_political\_opinions.csv` contains 149 rows: 40+99+10. The linked prompt dataset at Hugging Face commit `6f21c7e29e5acb3fdaa26c3161779f1d9e6ef2cb` [E4] contains 10,314 rows, including 2,980 unforced and 5,960 viewpoint prompts. Independent CSV counting, without executing author code, matches 44+48*20+30+149*20+149*40+17*10+17*10=10,314. The stated 150-item design gives 10,374 instead. The discrepancy is one original expanded into 60 prompts.\par
Alternative readings and impact. A different unpublished data version or an unreported exclusion could explain the mismatch, but the paper’s own 10,314 denominator corroborates the released counts. This is a Minor dataset-description inconsistency. The release does not establish which exact data generated every table, so neither the published model ordering nor a numerical effect on model scores is claimed to have changed. A documented exclusion/version map or correction to counts is needed; no experimental rerun was performed.\par
\section*{5. Other checked evidence and unresolved questions}
Table 2 aggregation remains Needs Evidence. Several percentages are compatible with base-question denominators rather than expanded prompt counts; that alone cannot distinguish majority aggregation from repeated identical labels or selection of one format. The Llama viewpoint row and the example on p41 warrant reconciliation, but source generations and aggregation code are absent. Do not infer a wrong table merely from a denominator pattern. The R1-Distill Universal Rights row totals 100.45\%, which cannot arise solely from ordinary rounding of an exhaustive six-category distribution. This is preserved as a local reporting observation, without a core-ranking claim.\par
Other bounded observations were not promoted to substantive findings. Appendix D.2’s 14/10,314 is 0.135738\%, not 0.001\%. The same paragraph prints a </text> terminator, whereas released `src/eval\_templates.py` lines 42-49 uses </think>; this is counterevidence to an implemented wrong-token error. Filename mismatches in the release are documented in `code\_review.md`, with no demonstrated link to published runs.\par
Rejected or clarified concerns. Appendix D.1 (p22) expressly explains choosing Avoidance for bad-faith conspiracies despite the flowchart’s Refusal branch; the data labels follow that documented choice. The evaluator’s single dominant final category matches the displayed prompt. Non-refusal is not unsafe compliance: a response may avoid the harmful request while offering safe related information, as the p60 example illustrates. The paper’s normative alternative-views section and trade-offs are not a mathematical impossibility theorem under an externally imposed neutrality metric.\par
Version and scope. The retrieved arXiv record lists v2 dated 3 June 2025; its pp17 and 23 retain the checked formal criterion and validation row. No correction was identified in the inspected publisher/arXiv records. OpenReview history was blocked by browser verification, so no exhaustive erratum claim is made. Original generation records, GPT-4o labels, human labels and aggregation scripts are missing from the pinned release; README output links are empty. New model calls would not reconstruct those historical outputs.\par
\section*{6. Conclusion and repair tracking}
The supported overall impact is Major. Problem 2 invalidates the advertised avoidance criterion within the core eight-technique formal framework. Problem 1 remains a Minor inconsistency in one human-validation summary, and Problem 3 remains a Minor dataset-description mismatch. Consequences for the model comparisons require additional source evidence; their survival does not reduce the severity of the failed formal criterion. The remaining contextual framework and normative recommendations have not been refuted by these checks. Human scientific decisions remain pending; author adoption has not been requested.\par
Item | Judgment | Repair verification | Remaining work\par
P1 validation row | Minor; mathematically inconsistent | No replacement validated | Aligned labels and statistic definition\par
P2 avoidance criterion | Major; core formalization admits a direct answer | Only set-distance implication checked | Specify complete operational rule\par
P3 sample counts | Minor; 149 vs 150 | Release counts independently verified | Historical data/exclusion map\par
Aggregation / rankings | Needs Evidence | Not repaired or reproduced | Generations, labels, analysis script\par
Normative framework | No demonstrated defect in checked argument | Not applicable | Future empirical evaluation as proposed\par
\section*{Supporting checks}
[E1] Fisher et al. (2025). Position: Political Neutrality in AI Is Impossible - But Here Is How to Approximate It. PMLR 267:81315-81374. https://proceedings.mlr.press/v267/fisher25a.html.\par
Fisher et al. arXiv:2503.05728v2 (3 June 2025). https://arxiv.org/abs/2503.05728v2. Compared only identified relevant passages; exact PDF retained.\par
Author repository: https://github.com/jfisher52/Approximation\_Political\_Neutrality. Pinned commit `e266bcf69689ea854897882b48fdf36f5679f239` (29 May 2025). Full acquired source tree and file hashes are recorded in code\_acquisition\_manifest.json and code\_repository\_tree.json.\par
Author prompt dataset: https://huggingface.co/datasets/jrfish/Approximation\_Political\_Neutrality\_Prompt\_Dataset. Pinned commit `6f21c7e29e5acb3fdaa26c3161779f1d9e6ef2cb`. `test\_data.csv` SHA-256: `f92688dd999b92e746e2546190e5b73993410442f8ac1bf125ccb7e09db58173`.\par
scikit-learn, cohen\_kappa\_score documentation, retrieved 24 October 2026. https://scikit-learn.org/stable/modules/generated/sklearn.metrics.cohen\_kappa\_score.html. Used for the standard kappa definition; no package execution.\par
\section*{Supporting checks}
The unchanged mathematical checks were reused by matching their source, script and result hashes; they were not rerun for this acceptance revision. This uses exact Fraction arithmetic and finite indicator/metric checks; `independent\_calculation\_results.json` contains the actual outputs. CSV counts and static source locators are in `code\_data\_counts.json`, `code\_hf\_acquisition.json` and `code\_static\_results.json`. Acquisition scripts download evidence only. `code\_static\_audit.py` parses source text/AST without importing author modules.\par
Source extraction used installed Poppler pdftotext -layout. Source pages 17 and 23 were rendered and visually checked.\par
No author code was executed; no model inference, training, paid API calls or historical empirical reproduction was performed. No human-review declaration or author-adoption claim was made.\par
This revision preserves the original inputs and changes only the documented severity/dependency reasoning and line wrapping. Human scientific decisions remain pending.\par
\end{document}
